\documentclass[11pt,a4paper]{ctexart}
\usepackage[margin=2.4cm]{geometry}
\usepackage{enumitem}
\usepackage{array}
\usepackage{booktabs}
\usepackage{amssymb}
\usepackage{xcolor}
\usepackage{titlesec}
\setlist{nosep,leftmargin=1.6em}
\titleformat{\section}{\large\bfseries}{\thesection}{0.6em}{}
\pagestyle{plain}

\begin{document}

\begin{center}
{\Large\bfseries AI 工具使用详情说明}\\[6pt]
{\normalsize 2026 年亚太地区大学生数学建模竞赛（APMCM）\quad A 题\quad 参赛队号：apmcm26300025}\\[4pt]
{\small 提交日期：2026 年 11 月 29 日}
\end{center}

\vspace{4pt}
\hrule
\vspace{10pt}

\section*{一、所用 AI 工具名称和版本}

\begin{table}[h]
\centering
\small
\begin{tabular}{@{}p{2.9cm}p{3.2cm}p{5.2cm}p{2.2cm}@{}}
\toprule
工具名称 & 版本 / 模型 & 提供方 & 使用日期 \\
\midrule
DeepSeek & DeepSeek-V4-Flash & 杭州深度求索人工智能基础技术研究有限公司 & 2026-11-29 \\
通义千问（Qwen） & Qwen3.8-Max & 阿里云计算有限公司 & 2026-11-29 \\
豆包 & 2.0 Pro & 北京抖音信息服务有限公司 & 2026-11-29 \\
\bottomrule
\end{tabular}
\end{table}

\section*{二、具体使用目的和应用环节}

\begin{enumerate}[label=\arabic*.]
\item \textbf{题目解析与问题分解}：辅助梳理五问的数据依赖、核对统一量 $F=W_1\setminus W_0$ 的可行性；统一量与各问输出由队伍确定。
\item \textbf{数据处理}：编写并调试遥感预处理脚本，包括 WGS-84 像元面积换算、分块秩归一化、Lee 滤波、Sentinel-1 三期影像互相关配准、WorldCover / OSM / WorldPop 叠加与分块 Otsu 阈值检验。
\item \textbf{模型构建}：在队伍确定的技术路线下，编写并调试双分支门控 SAR--光学 U-Net（含模态丢弃与有效性掩膜）、面积一致性损失、等渗校准与面积保持阈值、分裂保形区间、高斯混合分级与蒙特卡洛等级概率、熵--CRITIC--AHP--TOPSIS 单元评估、洪水切断路网上的选址--分配 MILP 与 $\varepsilon$-约束帕累托前沿的实现代码；网络结构与参数由队伍选定。
\item \textbf{求解与敏感性分析}：协助调试 HiGHS 求解器与 19 套方案的约束审计、Dirichlet 扰动、分块 bootstrap、地形护栏等敏感性脚本；方案设定与结果判读由队伍完成。
\item \textbf{可视化}：生成 21 张数据图、1 张场景示意图、技术路线图与融合架构图（矢量 PDF）、1 张 TikZ 集合逻辑图，并统一字体、配色与单页版式。
\item \textbf{文献检索与格式化}：通过 Crossref 核验 DOI 并生成 BibTeX 条目，剔除主题不符与重复的检索结果，统一参考文献格式。
\item \textbf{英文写作与语言润色}：起草各章英文正文、摘要与建议信，进行术语统一与语句润色。
\item \textbf{排版与编译检查}：XeLaTeX 编译、参考文献交叉引用核对、页眉与公式宽度修正、占位符清理。
\end{enumerate}

\section*{三、关键交互记录（重要提示词与回复摘要）}

\begin{enumerate}[label=\arabic*.]
\item \textbf{提示词}：“先核验官方附件标签的语义，再决定 Q1 识别器学什么。”
\textbf{回复摘要}：对全部 65.5 M 个有效像元统计后确认官方标签编码的是“新增淹没”，仅 2.37\% 的洪涝标签在灾前已为水体；队伍据此确认识别器直接学习 $F$，并在 Q3 强制 $F\cap W_0=\varnothing$。

\item \textbf{提示词}：“CBC 在 Windows 上求解敏感性时长时间无响应，换求解器重跑第五问。”
\textbf{回复摘要}：改用 HiGHS 后单次求解约 51 s，19 套方案全部通过约束审计（C10--C15 违规数均为 0），并给出 6 点帕累托前沿与 VoI$_1$=0.039、VoI$_2$=5.54 min。

\item \textbf{提示词}：“Sentinel-1 三期影像疑似存在整体偏移，先测偏移量再出图。”
\textbf{回复摘要}：测得南北向偏移 5.4--6.2 km（三期间互差 0.8 km），互相关配准后永久水体与季节性水体的暗像元比例明显上升，该修正写入正文配准与不确定性章节。

\item \textbf{提示词}：“把技术路线与融合架构画成论文插图，要求单页、无白边、文字矢量。”
\textbf{回复摘要}：输出 \texttt{fig\_roadmap.pdf} 与 \texttt{fig\_arch\_fusion.pdf}，数学符号正常排版，二者分别置于引言与第二问。

\item \textbf{提示词}：“论文写完后检查编译、页眉折行、公式出界与占位符。”
\textbf{回复摘要}：XeLaTeX + BibTeX 编译通过（0 错误、0 未定义引用）；将 10 个节标题改为短标题以消除页眉重叠，3 处超宽公式改为小号排版，删除首页保证书并填入队号。

\item \textbf{提示词}：“帮我核对论文里的数字与结果文件是否一致。”
\textbf{回复摘要}：逐项比对 $\tau_A$ 五折取值、校准面积 933.111/933.112 km$^2$、保形覆盖率 87.8\%、分级阈值 (0.0593, 0.5298, 0.9689)、硬瓦片 IoU 0.266$\to$0.562 等，均与结果文件一致。
\end{enumerate}

\section*{四、采纳内容和人工修改情况}

\textbf{采纳的 AI 输出}
\begin{itemize}
\item 五问预处理、特征与训练代码；队伍定结构后的门控融合与面积校准代码；阈值拟合与等级概率脚本；指标实现与权重计算；MILP 与帕累托求解脚本。
\item 全部 25 件图（21 张数据图、1 张场景图、2 张流程/架构图、1 张 TikZ 图）及其数据来源表。
\item 英文章节初稿、摘要、建议信初稿，以及参考文献的 DOI 核验与 BibTeX 条目。
\end{itemize}

\textbf{人工审核与修改}
\begin{itemize}
\item \textbf{数值核验}：逐项比对论文数字与结果文件，确认面积、覆盖率、阈值、分级计数、排名与救援指标一致；对校准面积相对误差等无法从舍入值反算的表述做了收紧。
\item \textbf{文献清理}：Crossref 检索结果中存在 1 条主题完全不符、3 条同一文献重复、以及 HTML 实体 \texttt{\&amp;} 与月份缩写不规范等问题，均人工剔除或改正，最终保留 20 条（含 3 条 AI 工具条目，即论文中上标 [1]--[3]）。
\item \textbf{结论表述}：第三问“净增为负”补“观测退水主导”解释；第五问“57 个单元不可达”归为情景约束。
\item \textbf{版式与合规}：删除首页保证书，填入参赛队号 apmcm26300025，改写页眉与超宽公式，补齐建议信的发文单位与日期，并按赛事要求另行提交本说明文件；正文已在摘要与第 1--10 节相应位置以上标 [1]--[3] 标注，使用声明见参考文献之后的《Report on Use of AI》。
\item \textbf{责任声明}：AI 工具用于辅助分析、编码、绘图与语言润色，全部模型设定、参数选择、结果解释与最终文字由参赛队审核确认，文责由参赛队承担。
\end{itemize}

\vspace{8pt}
\hrule
\vspace{6pt}
{\small 说明：本文件按赛事要求以 PDF 格式单独提交，命名遵循赛事《人工智能工具使用规定》第七条。}

\end{document}
